\section{Introduction}
An inference optimization can preserve the mathematical model and still change the behavior of a deployed numerical implementation. Quantized weights, finite-precision activations, physical cache extents and backend placement jointly determine the computation that produces a response. When a deployment's regression tests or downstream state are tied to an established numerical trace, replacing that trace with a different deterministic one does not satisfy the compatibility requirement. The practical question is therefore: \emph{which computations can be batched while preserving the existing reference?}

This question arises in both prompt processing and speculative decoding. Prompt tokens are known together; speculative decoding supplies several proposed continuation tokens for parallel verification~\cite{leviathan2023,chen2023}. Both expose independent query work inside a causal attention graph. Ordinary autoregressive continuation still generates one token at a time. The opportunity studied here is to retain the parallel work that is already available without changing how reference-sensitive operations are evaluated.

Existing approaches establish useful forms of determinism. Batch-invariant kernels stabilize per-token evaluation across batch sizes~\cite{thinkingmachines}. CoRun uses isolated prefill and fixed-shape decode~\cite{corun}. LLM-42 lets a fixed-shape verifier define a consistent execution, then repairs tokens and KV state~\cite{llm42}. Vosti connects canonical state to relational kernel contracts~\cite{vosti,vosticontract}. Our target is compatibility with a reference chosen before the optimization: the original quantized model bytes, numerical runtime and scalar execution recipe remain fixed.

Global serialization satisfies a simple execution constraint but removes batching from projections and feed-forward layers as well as attention. A more local approach evaluates attention separately for each query (PQ), leaving surrounding computation batched. Query-to-batch attention (QB) expresses this local organization through paired Q/mask views and shared K/V storage. However, agreement between PQ and QB is not sufficient: both can inherit a physical KV extent that differs from the scalar reference.

The decisive example is a four-query Llama verification crossing a padding boundary. The scalar executions use physical key extents $[768,768,768,1024]$, while the original batch uses $1024$ for every row. Identical visible values and zero masked probabilities do not recover the scalar result. Native operand interventions exchange the outcomes by changing physical extent. Grouping the first three queries at 768 and the last at 1024 repairs the discrepancy while retaining parallel work inside each group. Vulkan reproduces this mechanism and additionally requires its multirow Q8 projections to retain the reference's F32 activation path.

These observations lead to a method with three parts: extract each query's reference-sensitive execution requirements; partition adjacent queries into compatible groups; and lower the groups to existing attention operators while binding the selected policy to graph reuse. The method acts on physical execution requirements rather than input identity or a desired answer. Backend-specific arithmetic constraints accompany the common grouping rule. A separate Qwen3-14B result illustrates why placement also matters: the original scalar and batched paths place a head-normalization/RoPE boundary on different backends.

Our contributions are:
\begin{enumerate}[label=(\arabic*),itemsep=2pt]
\item A physical-prefix grouping rule and QB lowering that preserve scalar attention requirements while retaining surrounding batching and independent-query parallelism.
\item Native causal evidence on CUDA and Vulkan that distinguishes physical-prefix effects from activation quantization and backend placement, with targeted reversals and confirmation on previously unseen inputs.
\item Complete-request measurements that separate the cost of preserving the reference, the incremental value of QB over PQ, and the decision to enable drafting.
\end{enumerate}

The new Llama confirmation passes all 24 inputs on each backend. Vulkan's 4.95\% speculative QB/PQ reduction exceeds its registered engineering threshold; CUDA's 2.45\% reduction does not. On both Llama panels, target-only execution is cheaper than the corresponding speculative path. Earlier same-runtime scope comparisons show a complementary case: local repair makes drafting profitable for Qwen3 while SmolLM2 remains faster without it. The method preserves a useful execution choice; its cost must be measured in the intended request path.
